% Sample Applications, Chapter 4 -- Sales Commission (annex Sec 8.3).
% Values replayed through the shipped application (examples/
% commission.i23); @VLOOKUP as os/i123/core/eval.c fn_lookup computes it
% (the largest key not above the value; offset 1 = the second column).
% Figure screens/sa4.k123.
% GATES: test-i123-formula test-pubs-i123-examples test-pubs-screens
\Gates{test-i123-formula, test-pubs-i123-examples, test-pubs-screens}
\chapter{Sales Commission}

Each month the sales office computes the commission of each salesperson
from a table of rates: the more a salesperson sells, the higher the rate on
the whole amount.

\section{What You Enter}
\begin{itemize}
\item The rate table: the lowest sales of each bracket and its rate
\item Each salesperson's name and sales for the month
\end{itemize}

\section{Entering the Data}
\ix{"@VLOOKUP}
Type the rate table in A3..B6: the lowest sales of each bracket in
ascending order in column A, and the rate beside it in column B.  Type the
names in D3..D5 and the sales in E3..E5.  In F3 enter
\type{+e3*@vlookup(e3,\$a\$3..\$b\$6,1)} and copy it to F4 and F5.
\FlairShot[fig:sa4]{sa4}{Commissions from the rate table}

\begin{LotTable}{p{30mm}p{24mm}p{74mm}}
\LotHead{\textbf{Item} & \textbf{Address} & \textbf{Meaning}}
Rate table & A3..B6 & The lowest sales of each bracket and its rate\\
Name & D3..D5 & The salespeople\\
Sales & E3..E5 & Each one's sales for the month\\
Comm & F3..F5 & The commission: the sales times the rate of the bracket
  they fall in\LotEnd
\end{LotTable}

\section{Interpreting the Results}
\LotExample{examples/commission}
\fn{@VLOOKUP} finds the largest bracket in column A that is not more than
the sales, and returns the rate one column to its right.

\section{Changing the Template}
To change a rate, type it over the old one in column B.  To add a bracket
above the highest, type it in row 7, then change the range in the formulas
to \texttt{\$A\$3..\$B\$7} with \fkey{Edit}{F2}.  Column A must stay in
ascending order.
